\documentclass[12pt]{article}
\usepackage[margin=1in]{geometry}
\usepackage{amsmath,amssymb,amsthm}
\usepackage{algorithm,algorithmic}

\title{Problem 92: Square Digit Chains}
\author{Project Euler Solution}
\date{}

\begin{document}
\maketitle

\section{Problem Statement}

Define
\[
f(n)=\text{sum of the squares of the decimal digits of }n.
\]
Starting from a positive integer $n$, repeatedly apply $f$. We must count how many starting values $n<10^7$ eventually arrive at $89$.

\section{Reduction to finitely many states}

If $n<10^7$, then $n$ has at most $7$ decimal digits, so
\[
f(n)\le 7\cdot 9^2 = 567.
\]
Hence after one application of $f$, every chain lies in the finite set
\[
\{0,1,2,\dots,567\}.
\]
Therefore the eventual destination of $n$ depends only on $f(n)$.

By explicit computation on this finite set, the only eventual cycles are
\[
1\to 1
\]
and
\[
89\to145\to42\to20\to4\to16\to37\to58\to89.
\]
So every positive integer below $10^7$ ends at either $1$ or $89$.

\section{Counting by digit-square-sum}

Every integer $0\le n<10^7$ has a unique 7-digit decimal representation with leading zeros allowed. Thus we may count 7-digit strings
\[
d_1d_2\cdots d_7,\qquad d_i\in\{0,1,\dots,9\},
\]
by their square-digit sum
\[
d_1^2+\cdots+d_7^2.
\]

Let $C_k(s)$ denote the number of length-$k$ digit strings whose digit-square-sum equals $s$. Then
\[
C_0(0)=1,\qquad C_0(s)=0\ \text{for }s\ne 0,
\]
and the recurrence is
\[
C_{k+1}(s)=\sum_{d=0}^{9} C_k(s-d^2),
\]
where terms with negative index are treated as zero.

This follows because if the last digit is $d$, then the first $k$ digits must contribute exactly $s-d^2$.

After computing $C_7(s)$ for all $0\le s\le 567$, the answer is
\[
\sum_{\substack{0\le s\le 567\\ s\leadsto 89}} C_7(s),
\]
where $s\leadsto 89$ means the chain beginning at $s$ eventually reaches $89$.

The all-zero string contributes only to $s=0$, and $0$ does not reach $89$, so it does not affect the sum.

\section{Algorithm}

\begin{algorithm}[H]
\caption{Count starting values below $10^7$ that end at $89$}
\begin{algorithmic}[1]
\STATE For each $s=1,\dots,567$, determine whether iterating $f$ from $s$ ends at $89$
\STATE Initialize $\mathrm{dp}[0]=1$ and $\mathrm{dp}[s]=0$ for $s>0$
\FOR{$7$ digit positions}
    \STATE Create a new zero array $\mathrm{next}[0\ldots 567]$
    \FOR{each sum $s$}
        \FOR{each digit $d\in\{0,\dots,9\}$}
            \STATE $\mathrm{next}[s+d^2] \gets \mathrm{next}[s+d^2] + \mathrm{dp}[s]$
        \ENDFOR
    \ENDFOR
    \STATE $\mathrm{dp}\gets \mathrm{next}$
\ENDFOR
\STATE Sum $\mathrm{dp}[s]$ over all $s$ whose chain ends at $89$
\RETURN the sum
\end{algorithmic}
\end{algorithm}

\section{Complexity}

The DP uses
\[
7\times 568\times 10
\]
transitions and stores one array of length $568$. Hence
\[
\text{time}=O(7\cdot 568\cdot 10),\qquad \text{space}=O(568).
\]

\section{Answer}

\[
\boxed{8581146}
\]

\end{document}
